\section{Introduction}

Le raisonnement contrefactuel correspond à la capacité de se représenter ce qui aurait pu se passer si les choses avaient été différentes. Ce concept se trouve à l'intersection du domaine de la causalité et de l'explicabilité. En XAI (Explainable Artificial Intelligence), les explications contrefactuelles correspondent à proposer un changement minimal (selon différents critères possibles) qui permet de changer la sortie d'un système. Concernant la causalité, le raisonnement contrefactuel y tient une place centrale (notamment avec les travaux de Halpern et Pearl), on cherche à identifier les facteurs déterminants d'une situation et à évaluer leur contribution dans la situation finale.

Le travail proposé s'inscrit dans la continuité du cadre introduit dans \cite{sarmiento_formalising_2025} qui propose un formalisme général basé sur des systèmes à transitions d'état étiquetés (LTS) qui correspondent à des graphes orientés. Une trajectoire dans un LTS est entièrement déterminée par une politique déterministe qui associe à chaque état (sommet du graphe) un événement à exécuter (arête du graphe). L'objectif de ce travail est d'évaluer des contextes causaux contrefactuels, qui correspondent à des trajectoires alternatives à une trajectoire observée, selon deux axes : leur proximité à l'observation et leur vraisemblance. Sur la question de vraisemblance, on souhaite dans ce travail introduire une forme d'acceptabilité des contrefactuelles relative à des normes prescriptives ou statistiques, pour cela les arêtes des graphes sont augmentées par un attribut de coût relatif à des normes, fourni par un expert.

Pour répondre à cette problématique une fonction de coût modulaire est proposée, dont les composantes peuvent être adaptés en fonction de l'application et des contraintes d'un problème concret. Cette modularité s'appuie notamment sur des travaux en sciences cognitives qui s'intéressent, elles aussi, au raisonnement contrefactuel. 

Ce rapport est organisé de la façon suivante : dans une première partie, nous reprenons et adaptons les définitions et hypothèses formelles du cadre proposé par \cite{sarmiento_formalising_2025}. Puis, nous introduisons la fonction de coût et sa modularité (choix des différentes composantes, opérateurs, évaluation normative...). Dans une seconde partie plus applicative, nous développons une pipeline pour des applications concrètes. La première étape correspond à l'introduction d'un parser PDDL vers LTS afin de justifier la possibilité de travail dans le cadre formel introduit. Puis, nous appliquons notre formule à un exemple concret qui permet d'illustrer les différences introduites en fonction  de la configuration du coût.

